\documentclass[11pt,a4paper]{article}

\usepackage[utf8]{inputenc}
\usepackage[T2A]{fontenc}
\usepackage[ukrainian,english]{babel}
\usepackage[top=22mm,bottom=22mm,left=22mm,right=22mm]{geometry}
\usepackage{hyperref}
\usepackage{fancyhdr}
\usepackage{parskip}
\usepackage{tcolorbox}
\usepackage{enumitem}
\usepackage{listings}
\usepackage{xcolor}

\tcbuselibrary{skins,breakable}

\hypersetup{
    pdftitle={Workshop 10 -- Subagents \& Agent Teams: Handout},
    pdfauthor={Vadym Bondarenko},
    colorlinks=true,
    linkcolor=blue,
    urlcolor=cyan,
}

\pagestyle{fancy}
\fancyhf{}
\fancyhead[L]{\small\textit{Workshop 10 -- Subagents \& Agent Teams}}
\fancyhead[R]{\small\thepage}
\fancyfoot[C]{\small\textit{vcdev.me \textperiodcentered{} 2026}}
\renewcommand{\headrulewidth}{0.4pt}

\lstdefinelanguage{yaml}{
    keywords={true,false,null},
    sensitive=false,
    comment=[l]{\#},
    morestring=[b]',
    morestring=[b]"
}

\lstset{
    basicstyle=\ttfamily\small,
    breaklines=true,
    frame=single,
    backgroundcolor=\color{gray!10},
    rulecolor=\color{gray!40},
    xleftmargin=8pt,
    xrightmargin=8pt,
    aboveskip=6pt,
    belowskip=6pt,
}

\title{%
  {\Huge\bfseries Workshop 10: Subagents \& Agent Teams}\\[0.6em]
  {\large Коли спавнити, як визначати, як збирати у команду --- handout}
}
\author{Vadym Bondarenko \\ \small\href{mailto:vadym.bondarenko@vcdev.me}{vadym.bondarenko@vcdev.me}}
\date{2026}

\begin{document}

\maketitle

\begin{abstract}
Пост-воркшоп референс. Не повторює слайди --- сухе зведення команд, frontmatter-полів, decision-tree, чек-листів. Усі факти посилаються на офіційні docs (\href{https://code.claude.com/docs/en/subagents}{code.claude.com}).
\end{abstract}

\tableofcontents
\newpage

\section{Що таке subagent}

Окрема Claude-сесія з власним контекстним вікном, system prompt-ом і tool-restrictions. Викликається з main-сесії через \texttt{Agent} tool (раніше \texttt{Task}, перейменовано у v2.1.63). Не бачить історію main-розмови. Повертає у main лише summary.

\subsection{Built-in subagents}

\begin{tabular}{lll}
\hline
\textbf{Agent} & \textbf{Model} & \textbf{Tools} \\
\hline
Explore & Haiku & Read-only \\
Plan & Inherit & Read-only \\
General-purpose & Inherit & Усі \\
statusline-setup & Sonnet & (обмежені) \\
Claude Code Guide & Haiku & (обмежені) \\
\hline
\end{tabular}

\textbf{Explore} приймає \texttt{thoroughness}: \texttt{quick} / \texttt{medium} / \texttt{very thorough} (Claude обирає сам з фрази запиту).

\textbf{Subagent не може спавнити іншого subagent-а.} \texttt{Plan} існує спеціально, щоб запобігти infinite-nesting у plan-mode.

\section{Розташування файлів}

\begin{tabular}{lll}
\hline
\textbf{Scope} & \textbf{Path} & \textbf{Priority} \\
\hline
Managed settings & \texttt{<managed>/.claude/agents/<name>.md} & 1 (highest) \\
\texttt{--agents} CLI flag & (JSON) & 2 \\
Project & \texttt{.claude/agents/<name>.md} & 3 \\
Personal & \texttt{\textasciitilde/.claude/agents/<name>.md} & 4 \\
Plugin & \texttt{<plugin>/agents/<name>.md} & 5 (lowest) \\
\hline
\end{tabular}

\textbf{Discovery:} project subagent-и знаходяться walk-up з working dir. \texttt{--add-dir} \emph{не} сканує subagent-и.

\section{Frontmatter cheatsheet}

\subsection{Обов'язкові поля}

\begin{lstlisting}[language=yaml]
---
name: code-reviewer            # kebab-case, унікальний
description: Use proactively after code changes  # коли делегувати
---
\end{lstlisting}

\subsection{Часті поля}

\begin{lstlisting}[language=yaml]
tools: Read, Grep, Glob, Bash    # allowlist; інакше успадковує усі
disallowedTools: Write, Edit     # denylist (першим)
model: sonnet                    # sonnet | opus | haiku | full ID | inherit
permissionMode: dontAsk          # default | acceptEdits | auto | dontAsk | bypassPermissions | plan
maxTurns: 20                     # стелиш від running-amok
color: blue                      # red | blue | green | yellow | purple | orange | pink | cyan
\end{lstlisting}

\subsection{Розширені поля}

\begin{lstlisting}[language=yaml]
skills: [api-conventions, ...]   # preload skill-контенту (НЕ успадковує з parent)
mcpServers:                      # inline-MCP лише для цього subagent
  - playwright:
      type: stdio
      command: npx
      args: ["-y", "@playwright/mcp@latest"]
hooks:                           # PreToolUse / PostToolUse / Stop
  PreToolUse:
    - matcher: "Bash"
      hooks:
        - type: command
          command: "./scripts/validate.sh"
memory: project                  # user | project | local --- persistent dir
background: true                 # завжди у фоні
effort: high                     # low | medium | high | xhigh | max
isolation: worktree              # окрема git-копія репо
initialPrompt: "..."             # auto-submit при --agent main
\end{lstlisting}

\section{Тіло subagent-а}

Markdown = system prompt. \textbf{Замінює} default Claude Code system prompt. Тільки твоя інструкція + базові env-деталі (working dir).

Pattern:
\begin{enumerate}
  \item Один абзац: «You are X».
  \item Workflow з нумерованими кроками.
  \item Чек-ліст або decision-tree.
  \item Формат виводу.
\end{enumerate}

\section{Як викликати}

\subsection{Natural language}
\begin{lstlisting}
Use the test-runner subagent to fix failing tests
\end{lstlisting}

\subsection{@-mention (гарантовано)}
\begin{lstlisting}
@code-reviewer (agent) look at the auth changes
\end{lstlisting}

\subsection{Whole session as subagent}
\begin{lstlisting}
claude --agent code-reviewer
\end{lstlisting}

\subsection{Default for project}
\texttt{.claude/settings.json}:
\begin{lstlisting}
{ "agent": "code-reviewer" }
\end{lstlisting}

\section{Patterns}

\subsection{Parallel dispatch}

Multiple \texttt{Agent} tool-calls у одному assistant turn:

\begin{lstlisting}
> Use 3 Explore subagents IN PARALLEL to investigate auth, tests, build
\end{lstlisting}

Всі три виконуються concurrent. Wall-clock економить. Token-overhead не змінює (3$\times$ vs sequential).

\textbf{Caveat:} «Running many subagents that each return detailed results can consume significant context.» 3--5 normal, 10+ --- задумайся про agent team.

\subsection{Chain через файли}

Subagent-и не передають дані напряму. Передають через working directory (shared):

\begin{enumerate}
  \item Researcher: читає task.md, пише research.md
  \item Planner: читає research.md, пише plan.md
  \item Reviewer: читає обидва, пише review.md
  \item Main: читає review.md, презентує юзеру
\end{enumerate}

Кожен subagent повертає лише one-line confirmation. Main-context чистий.

\section{Ізольований контекст}

Subagent \textbf{отримує}:
\begin{itemize}
  \item Свій system prompt (тіло md)
  \item Task prompt від main (точно як написано)
  \item Базові env-деталі (working dir)
  \item Доступ до файлів (якщо tools дозволяють)
\end{itemize}

Subagent \textbf{не отримує}:
\begin{itemize}
  \item Історію main розмови
  \item Результати інших subagent-ів
  \item Skills, завантажені у parent (треба \texttt{skills:} у frontmatter)
\end{itemize}

\textbf{Правило:} task prompt має бути самодостатнім. Не «продовж те, що ми обговорювали».

\textbf{cd не персистить} між Bash-викликами всередині subagent-а і не впливає на main working dir.

\section{Foreground vs Background}

\begin{tabular}{lll}
\hline
& \textbf{Foreground} & \textbf{Background} \\
\hline
Блокує main? & Так & Ні \\
Permission prompts & Передаються тобі & Pre-approved при спавні \\
Clarifying questions & Передаються тобі & Tool-call падає \\
\hline
\end{tabular}

Перейти у фон вручну: «run this in the background» або \texttt{Ctrl+B}. Вимкнути всі: \texttt{CLAUDE\_CODE\_DISABLE\_BACKGROUND\_TASKS=1}.

\section{Resume subagent}

Кожен \texttt{Agent} call --- новий instance, чистий контекст. Щоб продовжити конкретного subagent-а:

\begin{itemize}
  \item Claude отримує agent ID після завершення
  \item Resumption через \texttt{SendMessage} tool з ID
  \item \texttt{SendMessage} \textbf{доступний лише при увімкнених agent teams}: \texttt{CLAUDE\_CODE\_EXPERIMENTAL\_AGENT\_TEAMS=1}
\end{itemize}

Транскрипти: \texttt{\textasciitilde/.claude/projects/\{project\}/\{sessionId\}/subagents/agent-\{id\}.jsonl}. Cleanup за \texttt{cleanupPeriodDays} (default 30).

\section{Обмеження subagent-ів}

\begin{enumerate}
  \item Не можуть спавнити інших subagent-ів (flat hierarchy).
  \item Немає mid-flight user interaction (background subagent не може попросити permission).
  \item Permissions фіксуються при спавні.
  \item Не успадковують skills з parent --- \texttt{skills:} у frontmatter явно.
  \item Plugin-subagent-и втрачають \texttt{hooks}, \texttt{mcpServers}, \texttt{permissionMode} (security).
  \item Forks (наслідування історії) --- лише interactive, не у \texttt{claude -p} headless.
  \item Working dir shared, але \texttt{cd} не персистить між Bash-викликами.
\end{enumerate}

\section{Agent Teams}

\textbf{Експериментально.} Увімкнути:
\begin{lstlisting}
{ "env": { "CLAUDE_CODE_EXPERIMENTAL_AGENT_TEAMS": "1" } }
\end{lstlisting}

Потребує Claude Code v2.1.32+.

\subsection{Subagents vs Agent teams}

\begin{tabular}{lll}
\hline
& \textbf{Subagents} & \textbf{Agent teams} \\
\hline
Контекст & Свій, summary назад до main & Свій, повністю незалежний \\
Комунікація & Лише з main & Teammate $\leftrightarrow$ teammate \\
Координація & Main керує & Shared task-list \\
Token cost & Менший & Значно більший \\
\hline
\end{tabular}

\subsection{Архітектура}

\begin{itemize}
  \item \textbf{Team lead} --- main session, що створила команду
  \item \textbf{Teammates} --- окремі Claude Code instances
  \item \textbf{Task list} --- shared список робіт; teammate-и клеймлять
  \item \textbf{Mailbox} --- inter-agent messaging (\texttt{SendMessage})
\end{itemize}

Storage: \texttt{\textasciitilde/.claude/teams/\{name\}/config.json}, \texttt{\textasciitilde/.claude/tasks/\{name\}/}.

\subsection{Display modes}

\begin{itemize}
  \item \texttt{in-process} (default fallback) --- усі teammate-и в main терміналі, Shift+Down цикл
  \item \texttt{tmux} --- кожен у своєму pane (потребує tmux або iTerm2 + \texttt{it2} CLI)
  \item Default \texttt{auto} --- split якщо вже у tmux, інакше in-process
\end{itemize}

Не працює у VS Code terminal, Windows Terminal, Ghostty (split mode).

\subsection{Reuse subagent definitions}

Lead може посилатися на існуюче subagent-визначення:
\begin{lstlisting}
> Spawn a teammate using the security-reviewer agent type to audit src/auth/
\end{lstlisting}

Teammate успадковує \texttt{tools} і \texttt{model}; тіло \textbf{дописується} до системного prompt-а (не замінює).

\textbf{Не успадковуються:} \texttt{skills}, \texttt{mcpServers} --- teammate вантажить їх з project + user settings нормально.

\subsection{Limitations agent teams}

\begin{itemize}
  \item Експериментально, поведінка може змінитися
  \item \texttt{/resume} не відновлює in-process teammate-ів
  \item Task status може лагати
  \item Shutdown повільний
  \item Одна team на сесію (no nested)
  \item Lead фіксований на час життя team
  \item Permissions фіксуються при спавні
\end{itemize}

\section{Decision tree}

\begin{lstlisting}
Задача переді мною
 ├─ Швидкий fix у відомих файлах? → Main conversation
 ├─ Verbose output (логи, search)? → Explore subagent
 ├─ Multi-step research + дія? → general-purpose subagent
 ├─ 3+ незалежних investigation-ів? → Parallel subagent dispatch
 ├─ research → plan → implement → review? → Subagent chain (file-based)
 ├─ Workers мають дискутувати? → Agent teams (experimental)
 └─ Те саме + повний контекст? → Fork (experimental, interactive only)
\end{lstlisting}

\textbf{Heuristic:} старт з main. Subagent --- коли verbose output загрожує main-context-у. Team --- коли координація між workers і є цінність.

\section{Production-чек-ліст}

Перед коммітом subagent-а у \texttt{.claude/agents/}:

\begin{itemize}[leftmargin=2em]
\item[$\square$] \texttt{name} kebab-case, унікальний
\item[$\square$] \texttt{description} front-load + «use proactively»
\item[$\square$] \texttt{tools} мінімально-достатній allowlist
\item[$\square$] \texttt{model} обираєш свідомо: \texttt{haiku} для простого, \texttt{sonnet} для аналізу
\item[$\square$] System prompt --- workflow з нумерованими кроками + формат виводу
\item[$\square$] Self-contained task prompt
\item[$\square$] Test взаємодії: і \texttt{@-mention}, і natural-language
\item[$\square$] \texttt{README.md} у \texttt{.claude/agents/} для команди
\end{itemize}

\section{Команди для вправ}

\subsection{Вправа 1 --- Explore}

\begin{lstlisting}
> Use the Explore subagent to find all places where DocRef is used
\end{lstlisting}

\subsection{Вправа 2 --- встановити custom subagent}

\begin{lstlisting}
mkdir -p ~/.claude/agents
cp starter/commit-message-writer.md ~/.claude/agents/
# перезапустити Claude Code
@commit-message-writer (agent)
\end{lstlisting}

\subsection{Вправа 3 --- паралель}

\begin{lstlisting}
> Use 3 Explore subagents IN PARALLEL to investigate auth, tests, build
\end{lstlisting}

\subsection{Вправа 4 --- chain}

\begin{lstlisting}
mkdir -p ~/.claude/agents workdir
cp starter/agents/*.md ~/.claude/agents/
cp starter/workdir/task.md workdir/

> Use chain-researcher to read workdir/task.md and write workdir/research.md.
  Then chain-planner reads research.md and writes plan.md.
  Then chain-reviewer reads both and writes review.md.
\end{lstlisting}

\section{Resources}

\begin{itemize}
\item \href{https://code.claude.com/docs/en/subagents}{code.claude.com/docs/en/subagents} --- Subagents reference
\item \href{https://code.claude.com/docs/en/agent-teams}{code.claude.com/docs/en/agent-teams} --- Agent Teams
\item \href{https://code.claude.com/docs/en/headless}{code.claude.com/docs/en/headless} --- Headless / Agent SDK
\item \href{https://code.claude.com/docs/en/cli-reference}{code.claude.com/docs/en/cli-reference} --- CLI flags
\item Цей репо: \texttt{workshops/10-subagents-and-teams/exercises/} і \texttt{solutions/}
\end{itemize}

\end{document}
